\section{Minor Examination, 2024--25 (Semester 1)}

\subsection*{Q1. Short definitions}

\textbf{Question.}
For each of the following terms, explain in one or two lines what they mean.

\begin{enumerate}
\item[(i)] \textbf{Principle of Isolation.}
Each process should run as if it has the machine to itself: its state and resources are protected so that one process cannot improperly observe or modify another process.

\item[(ii)] \textbf{TLB.}
A translation lookaside buffer (TLB) is a small, fast hardware cache of recently used page-table entries, mapping virtual page numbers to physical frame numbers and storing access permissions.

\item[(iii)] \textbf{Thrashing.}
Thrashing is the condition in which the system spends most of its time servicing page faults and moving pages between memory and disk, instead of executing useful work.

\item[(iv)] \textbf{Indirection (in the context of paging).}
The page table provides a level of indirection between virtual and physical
memory: a virtual page number indexes a PTE, which identifies the physical
frame containing that page.

\item[(v)] \textbf{Memory Hogs.}
Memory hogs are processes that consume a very large amount of memory, leaving
less physical memory available to other processes and potentially causing
heavy paging.

\item[(vi)] \textbf{Starvation.}
Starvation is indefinite postponement: a process remains ready or waiting but repeatedly fails to receive the CPU or another needed resource because other processes are continually preferred.

\item[(vii)] \textbf{Demand paging.}
Demand paging brings a virtual page into physical memory only when the process first references it; a reference to a non-resident page causes a page fault.
\end{enumerate}

\subsection*{Q2. Output of the \texttt{fork}/\texttt{exec}/\texttt{wait} program}

\textbf{Question.}
Predict the output of the following C program. Assume that function
\texttt{mywc()} takes a filename as parameter and writes the number of words
in it (integer) in standard output followed by filename. Justify your answer.

\begin{verbatim}
int main(int argc, char *argv[])
{
        printf("First: %d \n", (int) getpid());
        int rc = fork();
        if (rc == 0) {
                printf("Second: (pid:%d)\n", (int) getpid());
                char *myargs[3];
                myargs[0] = strdup("mywc");
                myargs[1] = strdup("abc.c");
                myargs[2] = NULL;
                execvp(myargs[0], myargs);
                printf("Second: %d"\n, (int) getpid());
        } else {
                wait(NULL);
                printf("First: %d\n Second:%d\n", rc, (int) getpid());
        }
return 0;
}
\end{verbatim}

\textbf{Solution.}
As printed, the program has no output because it does not compile: in the
second \texttt{printf} after \texttt{execvp}, the string closes before
\verb|\n|. The evidently intended line is
\begin{verbatim}
printf("Second: %d\n", (int) getpid());
\end{verbatim}
Under that intended correction, let the parent's PID be \(P\), the child's
PID be \(C\), and let \(W\) be the word count printed by \texttt{mywc}. If
\texttt{execvp()} succeeds, the output is
\begin{verbatim}
First: P
Second: (pid:C)
W abc.c
First: C
 Second:P
\end{verbatim}
Here \(P,C,W\) stand for their corresponding decimal integers. The first line
is printed before \texttt{fork()}. In the child, \texttt{fork()} returns zero,
so it prints its own PID \(C\), and successful \texttt{execvp()} replaces the
child's program image with \texttt{mywc}; consequently, the \texttt{printf()}
following \texttt{execvp()} is not executed. In the parent, \texttt{fork()}
returns the child's PID \(C\). The call to \texttt{wait(NULL)} ensures that the
parent's final two lines occur only after \texttt{mywc} has terminated; the
second argument of that final \texttt{printf()} is the parent's PID \(P\).

If \texttt{execvp()} fails, it returns to the child, so one extra line appears
before the parent can return from \texttt{wait()}:
\begin{verbatim}
First: P
Second: (pid:C)
Second: C
First: C
 Second:P
\end{verbatim}

\subsection*{Q3. Locality and the TLB}

\begin{enumerate}
\item[(i)] \textbf{Question: What do we mean by temporal and spacial locality
of a program? Explain with an example each.}

\textbf{Solution.}
Temporal locality means that a recently accessed address is likely to be
accessed again soon; repeatedly reading and updating a loop counter is an
example. Spatial locality means that addresses near a recently accessed
address are likely to be accessed soon; sequential traversal of an array is an
example.

\item[(ii)] \textbf{Question: How do temporal and spacial locality help in
making memory accesses faster in the context of TLB?}

\textbf{Solution.}
With temporal locality, repeated accesses commonly use the same virtual page
number (VPN), so its cached translation produces TLB hits. With spatial
locality, nearby addresses often lie on the same page and hence reuse one TLB
entry; even when traversal reaches nearby pages, their recently installed
entries may still be in the TLB. A TLB hit avoids an additional page-table
lookup in memory.

\item[(iii)] \textbf{Question: What do we mean by a fully associative TLB? Why
is that important?}

\textbf{Solution.}
In a fully associative TLB, a translation for any VPN may be stored in any TLB
slot, and the requested VPN is compared with all valid entries. Thus any
recently used page can occupy any free or selected slot, avoiding
set-conflict misses that a restrictive direct-mapped or set-associative
placement could cause in a small TLB.
\end{enumerate}

\subsection*{Q4. Page faults}

\textbf{Question.}
What is a page fault? Which part of modern computer handles page fault?
Describe the key steps that the component has to take while handling the page
fault.

\textbf{Solution.}
A page fault is a synchronous exception caused when a process accesses a
virtual address whose PTE says that the page is not present, or when the
requested access violates the PTE's protection bits. The MMU detects the
condition and raises the exception; the CPU traps into the operating-system
kernel's page-fault handler, which performs the actual handling.

For a legal reference to a non-resident page, the handler:
\begin{enumerate}
\item saves the process state and obtains the faulting virtual address and
access type;
\item checks that the address and requested access are legal (otherwise it
signals or terminates the process);
\item finds a free physical frame or selects a victim frame; if an evicted
page is dirty, writes it back to disk;
\item reads the required page from disk into the chosen frame;
\item updates the relevant PTE (and invalidates or updates stale TLB entries)
to record the frame, permissions, and present/valid bit; and
\item restarts the faulting instruction, which can now complete.
\end{enumerate}

\subsection*{Q5. Average memory-access time}

\textbf{Question.}
Find the average memory access time (AMAT) in the scenario, where the cache
miss rate is \(95\%\), time to access memory is \(50\) nanoseconds, and the
time to access the disk is \(20\)ms. How does AMAT change in this case if the
cache miss rate is reduced to \(1\%\)?

\textbf{Solution.}
The paper has a handwritten correction replacing the first ``miss'' by
``hit''; accordingly, the initial hit rate is \(95\%\) and the initial miss
rate is \(5\%\). Here memory is acting as a cache for disk. Using the OSTEP
weighted-average formula and
\(20\,\mathrm{ms}=20{,}000{,}000\,\mathrm{ns}\),
\[
\begin{array}{rcl}
\mathrm{AMAT}_{5\%\ \mathrm{miss}}
  &=& 0.95(50)+0.05(20{,}000{,}000)\ \mathrm{ns}\\
  &=&47.5+1{,}000{,}000\\
  &=&1{,}000{,}047.5\ \mathrm{ns}\\
  &\approx& 1.0000475\ \mathrm{ms}.
\end{array}
\]
When the miss rate is \(1\%\), the hit rate is \(99\%\), and hence
\[
\begin{array}{rcl}
\mathrm{AMAT}_{1\%\ \mathrm{miss}}
  &=&0.99(50)+0.01(20{,}000{,}000)\ \mathrm{ns}\\
  &=&49.5+200{,}000\\
  &=&200{,}049.5\ \mathrm{ns}\\
  &\approx&0.2000495\ \mathrm{ms}.
\end{array}
\]
Thus the AMAT falls from \(1{,}000{,}047.5\) ns to \(200{,}049.5\) ns, a
reduction by a factor of
\[
\frac{1{,}000{,}047.5}{200{,}049.5}\approx 4.999,
\]
i.e., approximately fivefold.

\subsection*{Q6. LRU versus FIFO}

\textbf{Question.}
Make an example with five page frames and eight pages where LRU performs
better than FIFO. What is one disadvantage of LRU over FIFO?

\textbf{Solution.}
Use pages \(\{0,1,\ldots,7\}\) and the reference string
\[
4,\ 6,\ 2,\ 1,\ 0,\ 4,\ 5,\ 3,\ 7,\ 4.
\]
Starting with five empty frames, the complete simulation is:

\begin{center}
\begin{tabular}{c|c|c|c|c|c}
Step & Ref. & FIFO frame contents & FIFO & LRU frame contents & LRU\\ \hline
1  & 4 & \(4,-,-,-,-\) & fault & \(4,-,-,-,-\) & fault\\
2  & 6 & \(4,6,-,-,-\) & fault & \(4,6,-,-,-\) & fault\\
3  & 2 & \(4,6,2,-,-\) & fault & \(4,6,2,-,-\) & fault\\
4  & 1 & \(4,6,2,1,-\) & fault & \(4,6,2,1,-\) & fault\\
5  & 0 & \(4,6,2,1,0\) & fault & \(4,6,2,1,0\) & fault\\
6  & 4 & \(4,6,2,1,0\) & hit   & \(4,6,2,1,0\) & hit\\
7  & 5 & \(5,6,2,1,0\) & fault & \(4,5,2,1,0\) & fault\\
8  & 3 & \(5,3,2,1,0\) & fault & \(4,5,3,1,0\) & fault\\
9  & 7 & \(5,3,7,1,0\) & fault & \(4,5,3,7,0\) & fault\\
10 & 4 & \(5,3,7,4,0\) & fault & \(4,5,3,7,0\) & hit\\ \hline
\multicolumn{2}{c|}{Total} & \multicolumn{2}{c|}{9 faults}
& \multicolumn{2}{c}{8 faults}
\end{tabular}
\end{center}

FIFO evicts page \(4\) at step 7 because it was loaded first, even though it
was just referenced at step 6; it therefore faults on page \(4\) at step 10.
LRU retains page \(4\), so it hits at step 10 and has fewer faults.

A disadvantage of exact LRU is its implementation overhead: the system must
track and update recency information (for example, timestamps or an ordered
stack) on memory references. FIFO needs only a simple queue recording page
arrival order.
